\section{技术浪潮切换：MLOps、RAG 与 Agentic AI}

\subsection{热度演化与真实需求}
在约 20 分钟处，讲者通过产业观察指出，MLOps、RAG、Prompt Engineering、Agentic AI 的热度迭代非常快。某些概念在一年内就从 ``新范式'' 走向 ``基础配置''。这并不意味着其失效，而意味着其被吸收进默认工程栈。

\begin{importantbox}{热度不等于价值，沉淀才等于价值}
RAG 的讨论热度下降，不代表 RAG 被淘汰；很多时候是因为它已成为系统默认能力，像数据库或缓存一样被隐式调用。
\end{importantbox}

\subsection{消费应用与企业应用的分化}
讲者把应用分为两类：
\begin{itemize}
    \item Consumer 场景：强调创意、娱乐、生产力，容忍一定不确定性；
    \item Enterprise 场景：强调正确性、稳定性、可追责，容忍错误的空间小得多。
\end{itemize}

这一区分决定了系统设计重点。消费场景常优先 ``体验增益''，企业场景必须先解决 ``可信交付''。同一模型在两类场景中，工程栈和治理要求差异极大。

\begin{knowledgebox}{同一模型，不同系统责任边界}
在消费产品里，系统更多是 ``加速器''；在企业系统里，系统是 ``安全阀''。前者可偏向试错，后者必须先构建保守默认、审计日志和回滚机制。
\end{knowledgebox}

\begin{warningbox}{错误迁移：把消费应用策略直接复制到企业场景}
许多团队在 PoC 阶段成功后直接推广到企业生产，忽视权限、合规、事实校验和 SLA 约束，最终导致 ``演示可行、生产不可用''。
\end{warningbox}

\subsection{阶段能力对比}
\begin{table}[H]
\centering
\begin{tabular}{p{2.8cm}p{3.6cm}p{3.6cm}p{3.2cm}}
\toprule
\textbf{阶段} & \textbf{核心能力} & \textbf{主要瓶颈} & \textbf{系统重点} \\
\midrule
MLOps 阶段 & 模型训练/部署流水线 & 实验管理与发布效率 & 平台化与自动化 \\
RAG 阶段 & 外部知识接入与问答增强 & 召回质量与延迟波动 & 检索链路与缓存策略 \\
Agent 阶段 & 规划、工具调用、长任务执行 & 长链路稳定性与可观测性 & 编排、流控、审计与恢复 \\
\bottomrule
\end{tabular}
\caption{从 MLOps 到 Agent 的系统责任扩张}
\end{table}

\subsection{本章小结}
浪潮切换的本质不是概念替换，而是能力沉淀与系统责任扩张。Agent 时代把系统设计推到前台。

